\documentclass[../main.tex]{subfiles}
\begin{document}

\chapter{Estructura del repositorio}
\label{anx:repositorio}

\section{Organización de ficheros}

El repositorio del proyecto se organiza de la siguiente manera:

\begin{lstlisting}[language=bash,
  caption={Árbol de directorios del repositorio del proyecto.},
  label={lst:repo}]
pfg_contenedores/
|-- README.md                     <- instrucciones de uso
|-- docker-compose.yml            <- definicion del stack completo
|-- .env.example                  <- plantilla de variables de entorno
|-- .gitignore
|-- generate_certs.sh             <- generacion de PKI con OpenSSL
|-- experiment_runner.py          <- orquestador de las 120 ejecuciones
|-- 00_ejecutar_todo.sh           <- script maestro de analisis
|-- 01_analisis_estadistico.py    <- Mann-Whitney U, Cliff delta, Fisher
|-- 02_figuras.py                 <- generacion de figuras PNG
|-- 03_tabla_latex.py             <- tabla LaTeX de resultados
|
|-- dockerfiles/
|   |-- fileserver.Dockerfile
|   |-- workstation.Dockerfile
|   |-- attacker.Dockerfile
|   `-- restic-sftp.Dockerfile
|
|-- lab-config/
|   |-- minio/
|   |   |-- restic-b-policy.json          <- politica IAM con permisos de borrado
|   |   |-- restic-c-policy.json          <- politica IAM sin permisos de borrado
|   |-- samba/
|   |   |-- smb.conf                      <- Conf Samba
|   |-- ssh-keys/                         <- claves SSH (privadas en .gitignore)
|   |   |-- lab_authorized_keys
|   |   |-- lab_runner_key
|   |   |-- lab_runner_key.pub
|   |   |-- restic_b_authorized_keys
|   |   |-- restic_b_key
|   |   |-- restic_b_key.pub
|   `--  wazuh/
|   |   |-- audit-rules-backup.rules
|   |   |-- audit-rules-victim.rules
|   |   |-- lab_ips                      <- IPs agentes
|   |   |-- local_rules.xml              <- 13 reglas custom del laboratorio
|   |   |-- local_rules.xml_auditd_old   <- reglas trabajo futuro auditd distinto kernel
|   |   |-- opensearch_dashboards.yml
|   |   |-- ossec_fileserver.conf        <- config agente fileserver
|   |   |-- ossec_restic-sftp.conf       <- config agente restic
|   |   `--  ossec_workstation.conf       <- config agente workstation
|
|-- scripts/
|   |-- reset_lab.sh              <- reset entre ejecuciones
|   |-- post_start.sh             <- post-arranque del stack
|   |-- verify_lab.sh             <- verificacion del estado del lab
|   |-- backup_volumes.sh         <- backup de volumenes Docker
|   |-- init_minio.sh             <- inicializacion de MinIO
|   |-- generate_test_files.sh    <- generacion de ficheros sinteticos
|   |-- fileserver-entrypoint.sh
|   |-- workstation-entrypoint.sh
|   |-- attacker-entrypoint.sh
|   |-- restic-sftp-entrypoint.sh
|   `-- playbooks/
|       |-- playbook_a.sh         <- estrategia A: sin backup
|       |-- playbook_b.sh         <- estrategia B: Restic+SFTP
|       `-- playbook_c.sh         <- estrategia C: MinIO Object Lock
|
|-- simulators/
|   |-- encryptor.py              <- Script encriptador dentro de rutas permitidas
|   `-- scenarios/
|       |-- e1_local.sh           <- T1486: cifrado AES-256-GCM
|       |-- e2_destroy.sh         <- T1490: destruccion de copias
|       |-- e3_lateral.sh         <- T1021.004: movimiento lateral SSH
|       `-- e4_exfil.sh           <- T1041: exfiltracion por SSH
|
|-- results/
|   |-- raw_20260724_174113.csv   <- datos del experimento completo
|   |-- resumen_celdas.csv        <- estadisticos por celda
|   |-- estadisticas.txt          <- analisis estadistico completo
|   |-- tabla_resultados.tex      <- tabla LaTeX para la memoria
|   |-- figura1_mttr_boxplot.png
|   |-- figura2_tasa_recuperacion.png
|   |-- figura3_tdetect.png
|   `-- figura4_mttr_scatter.png
|
`-- tfg_latex/                              <- memoria del PFG en LaTeX
    |-- acronimos.tex                       <- definicion de acronimos y glosario
    |-- bibliografia.bib                    <- referencias bibliograficas (BibTeX)
    |-- capitulos/
    |   |-- 01_introduccion.tex
    |   |-- 02_estado_arte.tex
    |   |-- 03_arquitectura.tex
    |   |-- 04_disenyo_experimental.tex
    |   |-- 05_resultados.tex
    |   |-- 06_discusion.tex
    |   |-- 07_conclusiones.tex
    |   |-- abstract.tex
    |   |-- agradecimientos.tex
    |   |-- anexo_a_reglas.tex
    |   |-- anexo_b_art.tex
    |   |-- anexo_c_repositorio.tex          <- este mismo anexo
    |   |-- referencias.tex
    |   `-- resumen.tex
    |-- imagenes/
    |   |-- figura1_mttr_boxplot.png
    |   |-- figura2_tasa_recuperacion.png
    |   |-- figura3_tdetect.png
    |   |-- figura4_mttr_scatter.png
    |   `-- rosadelosvientos.png
    |-- main.tex                            <- documento raiz (documentclass book)
    |-- Makefile                            <- automatizacion de la compilacion
    |-- portada.tex
    |-- README.md
    `-- referencias.bib
\end{lstlisting}

Cada uno de los directorios de primer nivel agrupa una responsabilidad
diferenciada del proyecto, siguiendo el principio de separar claramente el
código de infraestructura, el código de experimentación y la documentación
final. El directorio \texttt{dockerfiles/} contiene las definiciones de imagen
de los cuatro contenedores que requieren personalización más allá de una imagen
base oficial; los contenedores restantes (\texttt{wazuh-indexer},
\texttt{wazuh-manager}, \texttt{wazuh-dashboard} y \texttt{minio}) utilizan
directamente las imágenes oficiales publicadas por sus respectivos proyectos, sin
necesidad de un \texttt{Dockerfile} propio.

El directorio \texttt{lab-config/} centraliza toda la configuración estática
que los contenedores montan como volúmenes de solo lectura en tiempo de
arranque: las políticas \gls{iam} de MinIO que materializan la diferencia de
permisos entre las estrategias~B y~C descrita en la
sección~\ref{sec:almacenamiento-backups} del capítulo~\ref{C:arquitectura}, las
claves \gls{ssh} generadas específicamente para el laboratorio (excluidas del
control de versiones mediante \texttt{.gitignore} salvo las partes públicas), y
la configuración completa de Wazuh, incluyendo el fichero
\texttt{local\_rules.xml} reproducido íntegramente en el anexo~A.

Los directorios \texttt{scripts/} y \texttt{simulators/} separan
respectivamente la infraestructura de orquestación del experimento (reinicio del
laboratorio, verificación de estado, \emph{playbooks} de recuperación por
estrategia) de la lógica específica de cada escenario de ataque. Esta separación
permite que los cuatro escenarios descritos en la
sección~\ref{sec:variables-escenario} del capítulo~\ref{C:disenyo} puedan
extenderse o sustituirse sin modificar el orquestador central
\texttt{experiment\_runner.py}, en línea con el criterio de extensibilidad
modular reivindicado como contribución en el
capítulo~\ref{C:conclusiones}.

El directorio \texttt{results/} no se distribuye con el repositorio público
en su versión completa por motivos de tamaño, pero su estructura --- el CSV
crudo con las 120~filas del experimento, el resumen por celda y las cuatro
figuras generadas --- es la que se referencia directamente desde el
capítulo~\ref{C:resultados} de esta memoria. Finalmente, el directorio
\texttt{tfg\_latex/} contiene la memoria que el lector tiene actualmente
entre manos, organizada como un documento LaTeX modular mediante el paquete
\texttt{subfiles}: cada capítulo y anexo es un fichero independiente que puede
compilarse de forma aislada para revisión, además de integrarse en el documento
completo a través de \texttt{main.tex}.

\section{Reproducibilidad}

Para reproducir el experimento completo en un nuevo sistema con Docker
instalado se siguen estos pasos:

\begin{lstlisting}[language=bash,
  caption={Instrucciones completas de reproducción del experimento.},
  label={lst:reproduccion}]
# 1. Clonar el repositorio
git clone https://github.com/AlejandroMM2/Ransomware-simulator
cd pfg_contenedores

# 2. Requisito previo del kernel para OpenSearch
sudo sysctl -w vm.max_map_count=262144
echo "vm.max_map_count=262144" | sudo tee -a /etc/sysctl.conf

# 3. Copiar y editar el fichero de variables de entorno
cp .env.example .env
nano .env          # ajustar contrasenas si se desea

# 4. Generar certificados TLS para Wazuh
chmod +x generate_certs.sh
./generate_certs.sh

# 5. Construir las imagenes (primera vez: 15-20 min)
docker compose build

# 6. Arrancar el stack en dos etapas:

# Etapa A: arrancar el indexer y ejecutar securityadmin
docker compose up -d wazuh-indexer

# Esperar a que este sano
while [ "$(docker inspect -f '{{.State.Health.Status}}' \
  wazuh-indexer 2>/dev/null)" != "healthy" ]; do
    sleep 5; echo -n "."
done; echo " healthy"

# Inicializar plugin de seguridad de OpenSearch
docker exec wazuh-indexer bash -c '
  export OPENSEARCH_JAVA_HOME=/usr/share/wazuh-indexer/jdk
  bash /usr/share/wazuh-indexer/plugins/opensearch-security/tools/\
securityadmin.sh \
    -cd /usr/share/wazuh-indexer/opensearch-security/ \
    -nhnv \
    -cacert /usr/share/wazuh-indexer/certs/root-ca.pem \
    -cert /usr/share/wazuh-indexer/certs/admin.pem \
    -key /usr/share/wazuh-indexer/certs/admin-key.pem \
    -h localhost -p 9200 -icl
'

# Etapa B: arrancar el resto del stack
docker compose up -d

# 7. Verificar que todo esta operativo (esperar ~3 min)
bash scripts/post_start.sh
bash scripts/verify_lab.sh

# 8. Lanzar el experimento completo (~4-5 horas)
python3 experiment_runner.py

# 9. Generar analisis estadistico y figuras
bash 00_ejecutar_todo.sh
\end{lstlisting}

\section{Requisitos del sistema}

Para ejecutar el laboratorio se necesitan los siguientes recursos mínimos:

\begin{itemize}
  \item \textbf{CPU}: procesador x86\_64 de cuatro o más núcleos.
  \item \textbf{RAM}: 16\,GB (el stack en reposo consume $\approx$6\,GB).
  \item \textbf{Disco}: 30\,GB libres para imágenes Docker y volúmenes.
  \item \textbf{SO}: Linux con kernel $\geq$ 5.15 (no se requiere CachyOS).
    En distribuciones con kernel estándar el laboratorio puede
    complementarse con reglas de \texttt{auditd}.
  \item \textbf{Docker}: Engine~24.x o superior, Compose~v2.x.
  \item \textbf{Python}: 3.10 o superior con las bibliotecas
    \texttt{scipy}, \texttt{numpy}, \texttt{pandas} y
    \texttt{matplotlib} para el análisis estadístico.
\end{itemize}

\section{Licencia}

El código del laboratorio se publica bajo la licencia MIT, que permite
su uso, modificación y distribución con la única condición de mantener
el aviso de copyright original. La memoria del PFG se publica bajo la
licencia Creative Commons BY-NC-SA~4.0, que permite compartirla y
adaptarla con atribución, sin uso comercial y manteniendo la misma
licencia en los trabajos derivados.

\end{document}
